%--
\chapter{算法视角下的数学知识 (T1)}

%--
\begin{flushright}
\parbox{0.7\columnwidth}{\flushright \textit{\textsf{
    If you find the mathematics difficult, skip those early parts. 
    Later sections will be understandable 
    provided you are willing to forgo the deep insights mathematics 
    gives into the weaknesses of our current beliefs. 
    General results are always stated in words, 
    so the content will still be there but in a slightly diluted form.\\}}
\parbox{0.7\columnwidth}{\flushright \textsf{\textsf{\textemdash} 
    Richard W. Hamming \cite{Hamming05}}}
}
\end{flushright}


%----------------------------------------
\section*{授课提纲}

%----------
\subsubsection{I - 公理化视角}

\nid 从直觉开始构造
%--
\begin{itemize}
    \item 数学归纳法的基础：良序原理
    \item 自然数集合的公理化构造
    \item 数域的公理化构造：$N \Rightarrow Z \Rightarrow Q \Rightarrow R \Rightarrow C$
\end{itemize}

\nid 数学归纳法证明算法正确性
%--
\begin{itemize}
    \item 归纳法辨析：“所有马同色”谬误
    \item 归纳法应用：\alg{Euclid}，\alg{BubbleSort}，\alg{Prim}
\end{itemize}

\nid 认识和理解\ul{不变式}（invariant）的概念
%--
\begin{itemize}
    \item 归纳法中的“变”与“不变”
    \item 更大范围地应用不变式的概念
\end{itemize}

%----------
\subsubsection{II - 数学的概念，算法的视角}

\nid 基本数学概念
%--
\begin{itemize}
    \item 函数：$f(n)$
    \item 取整：$\floor{x}, \ceil{x}$
    \item 对数：$\log_a b$
    \item 阶乘：$n!$
\end{itemize}

\nid 基本级数
%--
\begin{itemize}
    \item 多项式级数：算术级数，平方和，任意$k$次方和（的渐近增长率）
    \item 几何级数：一般形式，2和$\frac{1}{2}$的特例，渐进增长率的速判
    \item 算术几何级数：“错项相减”法
    \item 调和级数
    \item 斐波那契数列
\end{itemize}

\nid 面向平均情况时间复杂度分析的期望值求解
%--
\begin{itemize}
    \item 指标随机变量
    \item 期望的线性性（linearality）
    \item 示例：排序算法平均情况时间复杂度的数学描述
\end{itemize}

%----------
\subsubsection{III - BF范式：分析驱动设计改进}

\nid BF范式（paradigm）
%--
\begin{itemize}
    \item BF是一个“相对性”的概念
    \item 基本模式：分析发现不足 $\Rightarrow$ 大O度量不足$\Rightarrow$不足推动改进
\end{itemize}


\nid 分析驱动改进
%--
\begin{itemize}
    \item “名人”问题；不变式：一次比较，确定淘汰一个
    \item 频繁项问题：$k=2$的特例；不变式：$|f| > |\neg f|$
\end{itemize}
